\documentclass[11pt]{article}

\usepackage[margin=1in]{geometry}
\usepackage[T1]{fontenc}
\usepackage[utf8]{inputenc}
\usepackage{lmodern}
\usepackage{booktabs}
\usepackage{array}
\usepackage{longtable}
\usepackage{tabularx}
\usepackage{listings}
\usepackage{xcolor}

\lstdefinestyle{cppstyle}{
  basicstyle=\ttfamily\small,
  keywordstyle=\color{blue},
  commentstyle=\color{teal},
  stringstyle=\color{magenta},
  showstringspaces=false,
  breaklines=true,
  frame=single
}

\title{DPS921 Lab Activity 3\\Sequential and CUDA Matrix Multiplication}
\author{Student Name: \underline{\hspace{2.4in}}\\Student ID: \underline{\hspace{2.7in}}\\Seneca Email: \underline{\hspace{2.35in}}}
\date{Submission Date: \underline{\hspace{2in}}}

\begin{document}

\maketitle

\section{Overview}
This report documents the implementation and performance analysis for Lab Activity 3. The lab compares a sequential CPU matrix multiplication program with two CUDA versions: a basic global-memory kernel and an optimized tiled shared-memory kernel. An optional Thrust-based implementation can also be included.

The matrix multiplication problem is defined as:
\[
C = A \times B, \quad c_{ij} = \sum_{k=1}^{N} a_{ik} b_{kj}
\]
All matrices are stored in row-major order using one-dimensional arrays.

\section{Environment}
\begin{tabularx}{\textwidth}{@{}lX@{}}
\toprule
Item & Details \\
\midrule
Operating system & \underline{\hspace{4.5in}} \\
Compiler for Question 1 & \underline{\hspace{4.2in}} \\
CUDA toolkit version & \underline{\hspace{4.1in}} \\
GPU model & \underline{\hspace{4.9in}} \\
Build commands & \small\texttt{g++ -O2 -std=c++17 question1.cpp -o question1} \\
& \small\texttt{nvcc -O2 question2.cu -o question2} \\
\bottomrule
\end{tabularx}

\section{Question 1: Sequential Matrix Multiplication}
\subsection{Implementation Summary}
The CPU version uses the standard three-loop matrix multiplication algorithm. For each output element \(C[i,j]\), the program computes the dot product of row \(i\) from matrix \(A\) and column \(j\) from matrix \(B\). The implementation accepts the matrix size as a command-line argument and reports the average execution time over five runs.

\subsection{Core CPU Algorithm}
\begin{lstlisting}[style=cppstyle,language=C++]
for (int i = 0; i < N; i++) {
    for (int j = 0; j < N; j++) {
        float sum = 0.0f;
        for (int k = 0; k < N; k++) {
            sum += A[i * N + k] * B[k * N + j];
        }
        C[i * N + j] = sum;
    }
}
\end{lstlisting}

\subsection{Timing Results}
Record at least five runs for each matrix size and report the average.

{\small
\setlength{\tabcolsep}{4pt}
\begin{longtable}{@{}>{\raggedright\arraybackslash}p{1.05in}*{5}{>{\centering\arraybackslash}p{0.78in}}>{\centering\arraybackslash}p{0.92in}@{}}
\toprule
Matrix size & R1 (ms) & R2 (ms) & R3 (ms) & R4 (ms) & R5 (ms) & Avg (ms) \\
\midrule
\endfirsthead
\toprule
Matrix size & R1 (ms) & R2 (ms) & R3 (ms) & R4 (ms) & R5 (ms) & Avg (ms) \\
\midrule
\endhead
256 $\times$ 256 & & & & & & \\
512 $\times$ 512 & & & & & & \\
1024 $\times$ 1024 & & & & & & \\
\bottomrule
\end{longtable}
}

\subsection{Correctness}
The CPU output was validated using the provided sample verification routine. Report the result here:

\begin{center}
\fbox{\parbox{0.85\textwidth}{Correctness result for Question 1: \underline{\hspace{2in}}}}
\end{center}

\section{Question 2: CUDA Matrix Multiplication}
\subsection{Part A: Basic CUDA Implementation}
The first CUDA version uses global memory only. Each thread computes exactly one output element of matrix \(C\). A two-dimensional block and grid configuration is used so that thread coordinates map naturally to matrix row and column positions.

\subsection{Basic CUDA Kernel}
\begin{lstlisting}[style=cppstyle,language=C++]
__global__ void matrixMultiplyBasicCUDA(float *A, float *B, float *C, int N)
{
    int row = blockIdx.y * blockDim.y + threadIdx.y;
    int col = blockIdx.x * blockDim.x + threadIdx.x;

    if (row < N && col < N) {
        float sum = 0.0f;
        for (int k = 0; k < N; k++) {
            sum += A[row * N + k] * B[k * N + col];
        }
        C[row * N + col] = sum;
    }
}
\end{lstlisting}

\subsection{Part B: Tiled Shared-Memory CUDA Implementation}
The optimized CUDA version uses shared memory tiling. A tile size of 16 $\times$ 16 was selected because it matches the block configuration required by the assignment and provides a practical balance between occupancy and shared-memory reuse on common GPUs.

Shared memory improves performance because threads in the same block can reuse data loaded from global memory. This reduces repeated global-memory reads and lowers overall memory access latency.

\subsection{Tiled CUDA Kernel Summary}
\begin{itemize}
\item Two shared tiles are declared: one for a sub-block of \(A\) and one for a sub-block of \(B\).
\item Each thread cooperatively loads one value from each matrix tile.
\item Threads synchronize before using the tile values.
\item Each thread accumulates a partial sum across all required tiles.
\item The final result is written to the corresponding output element in global memory.
\end{itemize}

\subsection{Question 2 Results}
Fill in the average timing results over five runs. Kernel time should include only the CUDA kernel execution measured using CUDA events. Total time should include allocation, host-to-device transfer, kernel launch, synchronization, device-to-host transfer, and cleanup.

{\small
\setlength{\tabcolsep}{4pt}
\begin{longtable}{@{}>{\raggedright\arraybackslash}p{1.15in}>{\centering\arraybackslash}p{0.95in}>{\centering\arraybackslash}p{1.15in}>{\centering\arraybackslash}p{1.1in}>{\centering\arraybackslash}p{0.95in}@{}}
\toprule
Matrix size & CPU avg & Basic kernel & Basic total & PASS/FAIL \\
\midrule
\endfirsthead
\toprule
Matrix size & CPU avg & Basic kernel & Basic total & PASS/FAIL \\
\midrule
\endhead
256 $\times$ 256 & & & & \\
512 $\times$ 512 & & & & \\
1024 $\times$ 1024 & & & & \\
\bottomrule
\end{longtable}
}

{\small
\setlength{\tabcolsep}{4pt}
\begin{longtable}{@{}>{\raggedright\arraybackslash}p{1.15in}>{\centering\arraybackslash}p{0.95in}>{\centering\arraybackslash}p{1.15in}>{\centering\arraybackslash}p{1.1in}>{\centering\arraybackslash}p{0.95in}@{}}
\toprule
Matrix size & CPU avg & Tiled kernel & Tiled total & PASS/FAIL \\
\midrule
\endfirsthead
\toprule
Matrix size & CPU avg & Tiled kernel & Tiled total & PASS/FAIL \\
\midrule
\endhead
256 $\times$ 256 & & & & \\
512 $\times$ 512 & & & & \\
1024 $\times$ 1024 & & & & \\
\bottomrule
\end{longtable}
}

\subsection{Performance Discussion}
Answer the required discussion points below.

\begin{enumerate}
\item Which GPU implementation performed better, the basic version or the tiled shared-memory version?
\item How much speedup did each GPU implementation achieve relative to the CPU implementation?
\item Why does shared memory improve performance for matrix multiplication?
\item What overheads are included in the total CUDA execution time but not in kernel execution time?
\end{enumerate}

\noindent Response:\\[0.3in]
\fbox{\parbox[t][2.0in][t]{\textwidth}{}}

\section{Question 3: Optional Thrust Component}
This section is optional. If completed, describe how manual device memory management was replaced with \texttt{thrust::host\_vector}, \texttt{thrust::device\_vector}, and \texttt{thrust::raw\_pointer\_cast()}.

\subsection{Thrust Results}
{\small
\setlength{\tabcolsep}{4pt}
\begin{longtable}{@{}>{\raggedright\arraybackslash}p{1.15in}>{\centering\arraybackslash}p{0.95in}>{\centering\arraybackslash}p{1.15in}>{\centering\arraybackslash}p{1.1in}>{\centering\arraybackslash}p{0.95in}@{}}
\toprule
Matrix size & CPU avg & Thrust kernel & Thrust total & PASS/FAIL \\
\midrule
\endfirsthead
\toprule
Matrix size & CPU avg & Thrust kernel & Thrust total & PASS/FAIL \\
\midrule
\endhead
256 $\times$ 256 & & & & \\
512 $\times$ 512 & & & & \\
1024 $\times$ 1024 & & & & \\
2048 $\times$ 2048 & & & & \\
\bottomrule
\end{longtable}
}

\subsection{Optional Discussion Prompts}
\begin{enumerate}
\item Which implementation performed best overall?
\item How much speedup did the best GPU version achieve over the CPU version?
\item What advantages does Thrust provide compared with manual CUDA memory management?
\item Under what conditions would you choose Thrust over handwritten CUDA memory management?
\end{enumerate}

\noindent Response:\\[0.3in]
\fbox{\parbox[t][2.2in][t]{\textwidth}{}}

\section{Conclusion}
Summarize the main findings of the lab in a short paragraph. State whether the program outputs were correct, whether the tiled kernel improved performance, and what was learned about the trade-off between kernel speed and total CUDA execution time.

\vspace{0.3in}
\fbox{\parbox[t][1.8in][t]{\textwidth}{}}

\end{document}
